\section{Требования к программе}

\subsection{Требования к функциональным характеристикам}

При проведении испытаний должны быть проверены требования к программе,
заданные в техническом задании и относящиеся к ее функциональным
характеристикам.

Испытуемая программа должна:
\begin{enumerate}
    \item обеспечивать генерацию таблиц набора данных TPC-H;
    \item формировать данные в соответствии с правилами и распределениями,
    принятыми в эталонном генераторе набора TPC-H;
    \item обеспечивать запуск из командной строки;
    \item принимать параметр масштаба генерации
    \texttt{--scale-factor};
    \item обеспечивать генерацию как полного набора таблиц, так и указанного
    оператором подмножества таблиц;
    \item обеспечивать запись результатов в формат вывода, предусмотренный
    конфигурацией конкретной сборки;
    \item обеспечивать задание каталога или URI вывода;
    \item выдавать справочную информацию и списки доступных
    таблиц и форматов;
    \item обеспечивать управление параметрами батчирования и записи выходных
    данных;
    \item обеспечивать корректную работу при больших значениях масштаба
    генерации, в том числе при значениях \texttt{scale factor}, превышающих
    100;
    \item диагностировать недопустимые значения параметров и
    завершаться с ненулевым кодом возврата при ошибке запуска.
\end{enumerate}

\subsection{Требования к надежности}

При недопустимых параметрах запуска, неподдерживаемом формате вывода,
некорректной стратегии записи, ошибке доступа к каталогу вывода либо иной
ошибке входных условий программа должна завершать выполнение управляемым
образом и выдавать оператору диагностическое сообщение, позволяющее установить
причину отказа.

При допустимых входных параметрах программа должна обеспечивать завершение
работы с формированием полного результата генерации либо завершение с
диагностикой при возникновении нештатной ситуации.

При работе на больших значениях масштаба генерации программа не должна
завершаться аварийно вследствие переполнения внутренних диапазонов генерации,
некорректной обработки ключевых значений или иных ошибок, связанных с ростом
объема синтезируемых данных.

\subsection{Требования к составу технических и программных средств}

Испытания должны проводиться на 64-разрядной операционной системе семейства
Linux при наличии собранного исполняемого файла \texttt{schengen\_main},
программных библиотек, необходимых для форматов хранения, включенных в текущую
сборку программы, а также прав доступа к целевому каталогу или URI вывода.
